% Local Self-Attention and Local Cross-Attention for the Phasor SSM —
% derivation document, companion to phasor_ssm_derivation.tex.
% Compile standalone: pdflatex local_attention_derivation.tex
%                     pdflatex local_attention_derivation.tex    (second pass for refs)

\documentclass[11pt]{article}
\usepackage{amsmath,amssymb,amsthm}
\usepackage{bm}
\usepackage{booktabs}
\usepackage{enumitem}
\usepackage[margin=1in]{geometry}
\usepackage{hyperref}

\newcommand{\R}{\mathbb{R}}
\newcommand{\C}{\mathbb{C}}
\newcommand{\im}{\mathrm{i}}
\newcommand{\dd}{\mathrm{d}}
\newcommand{\dt}{\Delta t}
\newcommand{\nh}{H}              % number of heads
\newcommand{\Dh}{D_{\mathrm{h}}} % per-head channel width
\newcommand{\sm}{\operatorname{sim}}
\newcommand{\eg}{e.g.}
\newcommand{\ie}{i.e.}

\theoremstyle{definition}
\newtheorem{definition}{Definition}
\newtheorem{proposition}{Proposition}
\newtheorem{remark}{Remark}

\title{Local Self- and Cross-Attention\\
for the Phasor State-Space Model}
\author{}
\date{}

\begin{document}
\maketitle

\begin{abstract}
\noindent
The Phasor SSM~\cite{phasor-ssm-derivation} represents every layer as a
diagonal complex state-space model and admits three equivalent
computational views: a recurrent discrete update, a continuous-time ODE,
and a causal convolution evaluated via FFT. The existing attention
primitives in the architecture (\texttt{PhasorAttention},
\texttt{SSMSelfAttention}, \texttt{SSMCrossAttention}) compute their
similarity score \emph{across time}, producing an $L\times L$ score
matrix that fuses information from every time slice at once. This
operation is intrinsically non-causal, so the attention layers cannot be
written as a per-slice discrete update, breaking the three-view
equivalence enjoyed by \texttt{PhasorDense} and \texttt{PhasorConv}.

We propose two attention primitives---\textbf{Local Self-Attention}
(LSA) and \textbf{Local Cross-Attention} (LCA)---that lift the
similarity computation off the time axis and onto a \emph{head} axis.
Each layer is defined on a single time slice and broadcast trivially
across time; the resulting layers inherit the three-view equivalence of
the surrounding SSM. We give the discrete, continuous, and parallel
forms, prove their equivalence, and relate the construction to local
attention variants in transformers~\cite{longformer,efficient-transformers}
and to modern Hopfield networks~\cite{hopfield-is-all-you-need}.
\end{abstract}

%======================================================================
\section{Motivation: the across-time score is non-causal}\label{sec:motivation}
%======================================================================

A Phasor SSM layer satisfies the per-channel equation
\begin{equation}\label{eq:ssm-base}
  \frac{\dd z_c}{\dd t} = k_c\,z_c(t) + \sum_j W_{cj}\,I_j(t),
  \qquad k_c = \lambda_c + \im\omega,
\end{equation}
where $\omega = 2\pi/T$ is the carrier shared by every channel of the
layer (the per-channel-$\omega$ rule, see~\cite{phasor-ssm-derivation},
\S3.5--3.7). After Dirac discretization with step $T$, the layer admits
the three equivalent forms
\begin{equation}\label{eq:ssm-three-views}
  z_c[n+1] = A_c\,z_c[n] + H_c[n], \qquad
  z_c = K_c * H_c, \qquad
  \frac{\dd z_c}{\dd t} = k_c\,z_c + W_c \cdot I(t),
\end{equation}
with $A_c = e^{k_c T}$, $B_c = (A_c-1)/k_c$, $K_c[n] = A_c^n$ for Dirac
inputs.

Vanilla self-attention, given a Phase tensor
$X\in[-1,1]^{C\times L\times B}$ and per-batch query/key/value
projections of the same shape, computes
\begin{equation}\label{eq:vanilla-attention}
  S_{l m}^{(b)} = \sm\bigl(Q_{:,l,b},\, K_{:,m,b}\bigr)
       \in \R^{L\times L\times B},
\end{equation}
\ie, a similarity between every pair of time slices. The score
$S_{l m}^{(b)}$ depends on $K_{:,m,b}$ for all $m$ including $m > l$, so
the operation is non-causal in time. Equivalently, the value mix
\(
Y_{:,l,b} = \sum_m V_{:,m,b}\,\bar S_{lm}^{(b)}
\)
is a contraction of $V$ along the time axis, which removes time as an
ambient dimension of the output.

These two properties have a single consequence: there is no per-slice
recurrence $z[l+1] = F(z[l], x[l])$ that reproduces vanilla attention.
Without such a recurrence, the layer cannot participate in the
discrete/continuous/parallel equivalence
of~\eqref{eq:ssm-three-views}, and so cannot be deployed on a spiking
substrate, nor as a streaming online operator, nor as a parallel
prefix scan.

Our remedy: replace the across-time similarity~\eqref{eq:vanilla-attention}
with an across-\emph{head} similarity. The score matrix becomes
$\R^{\nh\times \nh\times L\times B}$ (LSA) or
$\R^{A\times \nh\times L\times B}$ (LCA, with $A$ stored anchors), and
its computation factors slice-by-slice in $l$. The layer is then a
pointwise-in-$l$ map, trivially compatible with each of the three SSM
views.

\begin{remark}[Why this restriction is desirable]
Removing cross-time information flow from the attention block does not
remove cross-time computation from the model: the surrounding linear
SSM still propagates state through time via $A_c$, and the
HiPPO-initialized leakage spectrum supplies multi-timescale
memory~\cite{hippo}. Local attention therefore forces representation
learning to occur \emph{within} each phase code (\ie, in the
HD-VSA-typed channel space) rather than by re-routing previously stored
information across time slices. This is consistent with the Fourier
HRR algebra of~\cite{phasor-ssm-derivation}, \S4.
\end{remark}

%======================================================================
\section{Notation carried from the SSM derivation}\label{sec:notation}
%======================================================================

We use the symbols of~\cite{phasor-ssm-derivation} throughout:
\begin{itemize}[nosep,leftmargin=*]
  \item $C$ --- channel count of the surrounding SSM. Inside the
        head decomposition we write $D = C$ to match the conventional
        attention notation.
  \item $L$ --- sequence length, $B$ --- batch size.
  \item $\nh$ --- number of heads; per-head channel width
        $\Dh = D/\nh$, assumed integer.
  \item $A$ --- number of stored anchors (LCA only).
  \item $k_c = \lambda_c + \im\omega$ --- per-channel complex
        eigenvalue, with the carrier $\omega = 2\pi/T$ shared across
        channels of a given layer.
  \item Discrete kernels $A_c = e^{k_c T}$ (Dirac) or $e^{k_c\dt}$
        (ZOH), $B_c = (A_c-1)/k_c$, $K_c[n] = A_c^n\,(B_c)$.
  \item Phase variable $\theta_c\in[-1,1]$, related to the complex
        potential by $z_c = r_c\,e^{\im\pi\theta_c}$.
  \item $\sm(\bm\theta_1,\bm\theta_2) = \tfrac{1}{D}\sum_{c=1}^{D}
        \cos\!\bigl(\pi(\theta_{1,c} - \theta_{2,c})\bigr)$, the
        Fourier-HRR similarity of~\cite{phasor-ssm-derivation}, eq.~(22).
\end{itemize}

A Phase array of shape $(D, L, B)$ denotes
$\theta\in[-1,1]^{D\times L\times B}$ with $D$ the channel axis, $L$
the time axis, $B$ the batch axis. We write
$\theta_{:, l, b}\in[-1,1]^D$ for the slice at time $l$ and batch
index $b$.

%======================================================================
\section{Multi-head phase projections}\label{sec:multihead}
%======================================================================

Let $X\in[-1,1]^{D\times L\times B}$ be an input phase tensor. We
define query, key, and value tensors by three head-wise PhasorDense
projections (\cite{phasor-ssm-derivation}, \S3.2), each producing a
tensor of shape $(D, L, B)$:
\begin{equation}\label{eq:qkv-proj}
  Q = \mathrm{PhasorDense}_Q(X), \quad
  K = \mathrm{PhasorDense}_K(X), \quad
  V = \mathrm{PhasorDense}_V(X).
\end{equation}
The projections operate in Phase 3D (Dirac) mode: each output channel
evolves under the layer's per-channel $\lambda_c$ and the shared
carrier $\omega$.

We reshape the channel axis of $Q$, $K$, $V$ into $(\Dh, \nh)$:
\begin{equation}\label{eq:reshape-heads}
  Q,K,V\colon (D, L, B) \;\longmapsto\; (\Dh, \nh, L, B), \qquad
  D = \Dh\cdot \nh.
\end{equation}
The reshape is purely a view; it does not reorder phase data and does
not change the underlying carrier. We denote a single head's slice at
time $l$ as $Q_{:,h,l,b}\in[-1,1]^{\Dh}$.

\begin{remark}[Compatibility with the per-channel-$\omega$ rule]
The PhasorDense projections preserve the carrier $\omega$; the reshape
of~\eqref{eq:reshape-heads} only groups channels. Therefore every head
$h$ shares the same $\omega$, and head-wise Fourier-HRR similarity is
well-defined. The per-head leakage $\lambda_{:,h}$ may differ across
heads---this is a multi-timescale variation in line with HiPPO
initialization~\cite{hippo} and does not desynchronize the
representation.
\end{remark}

%======================================================================
\section{Local Self-Attention}\label{sec:lsa}
%======================================================================

\subsection{Per-slice definition}

\begin{definition}[Local Self-Attention]\label{def:lsa}
Fix a time slice $l$ and a batch index $b$. Let $Q,K,V$ be the
head-reshaped tensors of~\eqref{eq:reshape-heads} with shape
$(\Dh, \nh, L, B)$ and let $\beta\in\R$ be a scalar scale parameter.
\emph{Local Self-Attention} is the map
\[
  \mathrm{LSA}_{l,b}\colon
    (Q_{:,:,l,b},\,K_{:,:,l,b},\,V_{:,:,l,b})
    \longmapsto Y_{:,l,b}\in[-1,1]^D
\]
defined by
\begin{align}
  s_{h h'}^{(l,b)}
    &= \sm\!\bigl(Q_{:,h,l,b},\, K_{:,h',l,b}\bigr)
       \in [-1, 1],
       \label{eq:lsa-score} \\
  \tilde s_{h h'}^{(l,b)}
    &= \frac{1}{\nh}\,\exp\!\bigl(\beta\,s_{h h'}^{(l,b)}\bigr),
       \label{eq:lsa-scale} \\
  \tilde Y_{:,h,l,b}
    &= \sum_{h'=1}^{\nh} \tilde s_{h h'}^{(l,b)}
       \cdot e^{\im\pi V_{:,h',l,b}}
       \in \C^{\Dh},
       \label{eq:lsa-mix} \\
  Y_{:,l,b}
    &= \texttt{reshape}_{(\Dh,\nh)\to D}\!\bigl(
         \tfrac{1}{\pi}\arg \tilde Y_{:,:,l,b}\bigr)
       \in [-1,1]^D.
       \label{eq:lsa-reduce}
\end{align}
The map is broadcast independently over $l\in\{1,\dots,L\}$ and
$b\in\{1,\dots,B\}$ to produce a full output of shape $(D, L, B)$.
\end{definition}

The score~\eqref{eq:lsa-score} is the Fourier-HRR similarity
between two phase vectors of length $\Dh$; it lives in $[-1, 1]$ with
$+1$ at coincidence and $-1$ at antipode. Equation~\eqref{eq:lsa-scale}
applies the scaled exponential used by \texttt{score\_scale}
(\cite{phasor-ssm-derivation} companion code, \texttt{src/network.jl})
and yields a non-negative row-stochastic-up-to-normalization weighting
across the second head index $h'$. The mix~\eqref{eq:lsa-mix} is a
bundling (complex addition) in the sense of \cite{phasor-ssm-derivation},
\S4.2, weighted by the head-similarity scores. Equation~\eqref{eq:lsa-reduce}
extracts phase via $\arg$ and reshapes back to $(D, L, B)$.

\subsection{Three execution modes}

Because $\mathrm{LSA}_{l,b}$ depends only on slice $l$, the layer
composes with the surrounding SSM in all three views of~\eqref{eq:ssm-three-views}.
Let $f_\beta\colon[-1,1]^{D\times B}\to[-1,1]^{D\times B}$ denote the
slice-restricted LSA map at a fixed $\beta$.

\paragraph{Discrete (recurrent).}
With an enclosing Phasor SSM layer of per-channel eigenvalues $k_c$ and
discrete transition $A_c = e^{k_c T}$,
\begin{equation}\label{eq:lsa-discrete}
  z_c[l+1] = A_c\,z_c[l] + \sum_j W_{cj}\,e^{\im\pi\,f_\beta(x[l])_j},
\end{equation}
where $x[l]\in[-1,1]^{D\times B}$ is the input slice at time $l$.
This is identical to~\eqref{eq:ssm-three-views} with the input phase
replaced by $f_\beta(x[l])$, applied slice-by-slice.

\paragraph{Continuous (ODE).}
Let $\tilde X(t)\in[-1,1]^{D\times B}$ be a continuous-time
interpolation of the input phase slices. The continuous LSA-driven SSM
satisfies
\begin{equation}\label{eq:lsa-continuous}
  \frac{\dd z_c}{\dd t}
    = k_c\,z_c(t) + \sum_j W_{cj}\,I_j(t),
  \qquad I_j(t) = e^{\im\pi\,f_\beta(\tilde X(t))_j},
\end{equation}
which is the input-replacement version of~\eqref{eq:ssm-base}.

\paragraph{Parallel.}
Define the input-replaced sequence
\(
  \hat X[l] = f_\beta(x[l])
\)
for $l=0,\dots,L-1$ and form the Dirac encoding
$\hat H_c[l] = \sum_j W_{cj}\,e^{k_c\,\delta t_{j}[l]}$ on the slice
$\hat X[l]$. Then
\begin{equation}\label{eq:lsa-parallel}
  z_c = K_c * \hat H_c,
  \qquad K_c[n] = A_c^n,
\end{equation}
matching~\eqref{eq:ssm-three-views}.

\begin{proposition}[Three-view equivalence for LSA]\label{prop:lsa-equiv}
Let $(z^{(\mathrm{rec})}, z^{(\mathrm{conv})}, z^{(\mathrm{ode})})$
denote the outputs of~\eqref{eq:lsa-discrete}, \eqref{eq:lsa-parallel},
and~\eqref{eq:lsa-continuous} respectively, with the convention
$\tilde X(l\,T) = x[l]$ on the sample lattice. Then
\[
  z^{(\mathrm{rec})}[l] = z^{(\mathrm{conv})}[l]
  = z^{(\mathrm{ode})}(l\,T),
  \qquad l = 0, 1, \dots, L-1.
\]
\end{proposition}

\begin{proof}
The map $f_\beta$ acts pointwise in $l$, so the sequence
$\hat X[l] = f_\beta(x[l])$ is well-defined and the encoding
$\hat H_c[l]$ depends only on $\hat X[l]$ via the Dirac formula
(\cite{phasor-ssm-derivation} eq.~(13)). With this encoding,
\eqref{eq:lsa-discrete} and \eqref{eq:lsa-parallel} reduce exactly to
the Dirac discrete recurrence and the Dirac causal convolution of
\cite{phasor-ssm-derivation}, \S3.3 (\eqref{eq:ssm-three-views} above),
which are equivalent by that paper's Proposition on
view-equivalence (\eqref{eq:ssm-three-views} unrolls into the causal
convolution algebraically).

For the continuous--discrete equivalence, restrict the LSA-driven ODE
of~\eqref{eq:lsa-continuous} to a single period $[l T,(l+1)T)$. On this
interval, $\tilde X(t)$ takes the constant value $x[l]$ (by the
zero-order interpolation convention), so the forcing
$I_j(t) = e^{\im\pi f_\beta(x[l])_j}$ is the Dirac response of the same
slice, and integration over the period reproduces the recurrence
\eqref{eq:lsa-discrete} with $A_c = e^{k_c T}$ (\cite{phasor-ssm-derivation},
\S3.3, Dirac discretization).
\end{proof}

\subsection{Connection to Fourier HRR}

The LSA score~\eqref{eq:lsa-score} computes the cosine similarity of
two unit-modulus complex vectors of length $\Dh$; this is the
Fourier-HRR similarity of \cite{phasor-ssm-derivation}, eq.~(22),
restricted to a per-head subspace. The mix~\eqref{eq:lsa-mix} is the
bundling operation of \cite{phasor-ssm-derivation}, eq.~(18), weighted
by the exponentiated similarity scores. The reshape and $\arg$
of~\eqref{eq:lsa-reduce} reproduce the
phase-extraction nonlinearity of \cite{phasor-ssm-derivation},
\S3.5---\ie, LSA is a similarity-weighted bundle in HD-VSA space,
performed once per time slice.

\begin{proposition}[LSA as interference-weighted bundle]\label{prop:lsa-hd}
Let $s = (\tilde s_{h h'}^{(l,b)})_{h,h'=1}^{\nh}\in\R^{\nh\times \nh}$
be the row-normalized score matrix of~\eqref{eq:lsa-scale} and let
$V_h = e^{\im\pi V_{:,h,l,b}}\in\C^{\Dh}$. Then
\[
  \tilde Y_{:,h,l,b} = (V \cdot s^\top)_{:,h}
  = \texttt{bundle}_{h'}\!\bigl(s_{h h'}\cdot V_{h'}\bigr),
\]
where $\texttt{bundle}$ is the Fourier-HRR bundling of
\cite{phasor-ssm-derivation}, eq.~(18). Thus LSA is the
similarity-weighted superposition of head-restricted Fourier-HRR
vectors, exactly mirroring the temporal-bundling form of
\cite{phasor-ssm-derivation}, eq.~(21), but on the head axis.
\end{proposition}

\begin{proof}
Direct from~\eqref{eq:lsa-mix} and the definition of bundling as
complex addition.
\end{proof}

%======================================================================
\section{Local Cross-Attention}\label{sec:lca}
%======================================================================

\subsection{Per-slice definition with stored anchors}

\begin{definition}[Local Cross-Attention]\label{def:lca}
Let $Q^{*}\in[-1,1]^{D\times A}$ be a bank of \emph{anchor phase
vectors} (trainable, time-invariant), and let $X\in[-1,1]^{D\times L\times B}$
be the input. Reshape $Q^{*}$ to $(\Dh, \nh, A)$ and project $X$ to
$K, V\colon (\Dh, \nh, L, B)$ via PhasorDense as
in~\eqref{eq:qkv-proj}--\eqref{eq:reshape-heads}.
\emph{Local Cross-Attention} is the map
$(Q^{*}, X)\mapsto Y\in[-1,1]^{D\times L\times B}$ defined by
\begin{align}
  s_{a h}^{(l,b)}
    &= \sm\!\bigl(Q^{*}_{:,a,h},\, K_{:,h,l,b}\bigr)
       \in [-1,1],
       \label{eq:lca-score} \\
  \tilde s_{a h}^{(l,b)}
    &= \frac{1}{\nh}\exp\!\bigl(\beta\,s_{a h}^{(l,b)}\bigr),
       \label{eq:lca-scale} \\
  \tilde Y_{:,h,l,b}
    &= \sum_{a=1}^{A}\sum_{h'=1}^{\nh} \tilde s_{a h'}^{(l,b)}
       \cdot M_{a h h'}\cdot e^{\im\pi V_{:,h',l,b}}
       \in \C^{\Dh},
       \label{eq:lca-mix} \\
  Y_{:,l,b}
    &= \texttt{reshape}_{(\Dh,\nh)\to D}\!\bigl(
         \tfrac{1}{\pi}\arg \tilde Y_{:,:,l,b}\bigr),
       \label{eq:lca-reduce}
\end{align}
where $M\in\R^{A\times \nh\times \nh}$ is a trainable head-mixing
weight (default $M_{a h h'} = \delta_{h h'}/A$, in which case
\eqref{eq:lca-mix} collapses to $\tfrac{1}{A}\sum_a \tilde s_{a h}\,
e^{\im\pi V_{:,h,l,b}}$).
\end{definition}

The construction mirrors LSA but with the per-slice query replaced by
the fixed anchor bank. The score matrix has shape
$(A, \nh, L, B)$ instead of $(\nh, \nh, L, B)$; the cost per slice is
$O(A\,\nh\,\Dh)$, linear in every dimension.

\subsection{Three execution modes}

The LCA map is again pointwise in $l$. The three views are exactly as
for LSA, replacing $f_\beta(x[l])$ with $g_\beta(x[l]; Q^{*})$
throughout:
\begin{align}
  &\textbf{Discrete:} \quad
    z_c[l+1] = A_c\,z_c[l] + \sum_j W_{cj}\,
      e^{\im\pi\,g_\beta(x[l]; Q^{*})_j},
       \label{eq:lca-discrete} \\
  &\textbf{Continuous:} \quad
    \frac{\dd z_c}{\dd t} = k_c\,z_c(t) + \sum_j W_{cj}\,
      e^{\im\pi\,g_\beta(\tilde X(t); Q^{*})_j},
       \label{eq:lca-continuous} \\
  &\textbf{Parallel:} \quad
    z_c = K_c * \hat H_c,
    \quad \hat H_c[l] = \sum_j W_{cj}\,
      e^{k_c\,\delta t_j(g_\beta(x[l]; Q^{*}))}.
       \label{eq:lca-parallel}
\end{align}

\begin{proposition}[Three-view equivalence for LCA]\label{prop:lca-equiv}
The discrete, parallel, and (sampled) continuous outputs of LCA agree
at every lattice point $l T$, $l = 0,\dots,L-1$.
\end{proposition}

\begin{proof}
Identical to the proof of Proposition~\ref{prop:lsa-equiv}: the anchor
bank $Q^{*}$ is time-invariant so $g_\beta(\cdot; Q^{*})$ remains
pointwise in $l$, and the surrounding SSM is unchanged.
\end{proof}

\begin{remark}[Cost]
LCA's parallel-mode complexity is $O(A\,\nh\,\Dh\,L\,B)$, linear in
every dimension. The trainable anchor bank costs $A\,D$ parameters
(plus $A\,\nh^2$ for $M$ if used).
\end{remark}

\subsection{LCA as a modern Hopfield retrieval}

Substituting~\eqref{eq:lca-score} into~\eqref{eq:lca-scale} and
identifying $\beta$ as an inverse temperature yields the score
\[
  \tilde s_{a h}^{(l,b)}
    = \frac{1}{\nh}\exp\!\Bigl(
        \frac{\beta}{\Dh}\sum_{c=1}^{\Dh}
          \cos\!\bigl(\pi(Q^{*}_{c,a,h} - K_{c,h,l,b})\bigr)\Bigr).
\]
This is the modern Hopfield update of~\cite{hopfield-is-all-you-need}
restricted to the unit torus and using the Fourier-HRR inner product
$\tfrac{1}{\Dh}\langle Q^{*}_{:,a,h}, K_{:,h,l,b}\rangle$ in place of
the Euclidean inner product. When $V = K$ and $M = \delta_{h h'}/A$,
the mix~\eqref{eq:lca-mix} reproduces the Hopfield retrieval rule
$\xi^{\mathrm{new}} = X\,\mathrm{softmax}(\beta\,X^\top\xi)$ with
$X$ playing the role of stored patterns.

\begin{proposition}[Anchors as stored patterns]\label{prop:lca-hopfield}
Define the per-head pattern matrix $\Xi_h \in \C^{\Dh \times A}$ with
columns $\Xi_{:,a,h} = e^{\im\pi Q^{*}_{:,a,h}}$, and assume
$V = K$ and $M_{a h h'} = \delta_{h h'}/A$. Then
\[
  \tilde Y_{:,h,l,b}
  = \frac{1}{\nh A}\,\Xi_h \cdot \exp\!\bigl(\beta\,
       \mathrm{Re}\,\Xi_h^{\dagger}\,e^{\im\pi K_{:,h,l,b}} / \Dh\bigr),
\]
which is the modern-Hopfield retrieval rule
of~\cite{hopfield-is-all-you-need} with stored patterns $\Xi_h$ and
query $e^{\im\pi K_{:,h,l,b}}$, up to the per-head normalization
factor $1/(\nh A)$.
\end{proposition}

\begin{proof}
Substitute the Fourier-HRR similarity~\eqref{eq:lca-score} into
\eqref{eq:lca-mix} and collect terms. The real part of
$\Xi_h^{\dagger}\,e^{\im\pi K_{:,h,l,b}}$ is exactly $\Dh \cdot
\sm(Q^{*}_{:,a,h}, K_{:,h,l,b})$, so the exponential reproduces
\eqref{eq:lca-scale}. The contraction $\Xi_h \cdot (\cdot)$ then
realizes \eqref{eq:lca-mix} under $V=K$.
\end{proof}

Equivalently, LCA inherits modern Hopfield's exponential pattern
capacity ($A_{\max}\sim e^{c\Dh}$, see~\cite{hopfield-is-all-you-need},
Theorem 3), modulo the unit-torus restriction; we leave a precise
capacity bound as future work.

%======================================================================
\section{Comparison with related local-attention constructions}
\label{sec:related}
%======================================================================

\subsection{Window-restricted local attention}

Sliding-window transformer attention~\cite{longformer,longformer-followups}
reduces vanilla attention's $O(L^2)$ cost to $O(L\,w)$ by restricting
the time-axis softmax to a window of width $w$. The score
remains across-time and depends on $w$ future-and-past slices, so the
operation is non-causal except in the unidirectional variant. Our LSA
differs structurally: it eliminates the time axis from the score
entirely, replacing it with a head axis of size $\nh^2$. The per-slice
cost is $O(\nh^2\,\Dh)$, independent of $L$, and the layer is
\emph{strictly local in time}---it depends on no other time slice.
This is a stronger restriction than windowed attention and a stronger
guarantee for spiking deployment, where causal locality is a hardware
prerequisite.

\subsection{State-space attention}

Mamba's selective state-space model~\cite{mamba} makes the SSM
transition $A_c$ input-dependent. RetNet~\cite{retnet} chunks
retention into per-block recurrences. S5~\cite{s5} pairs a diagonal
SSM with a linear-attention output projection. All three blend SSM and
attention by making transition parameters depend on the input. Our
construction does the opposite: the SSM transition is fixed (constant
$k_c$ as in~\cite{phasor-ssm-derivation}), and the attention is
applied pointwise in time. The two design axes are orthogonal; LSA/LCA
could in principle be composed with selective-state mechanisms,
though we do not pursue that here.

\subsection{Modern Hopfield networks}

\cite{hopfield-is-all-you-need} established that a continuous Hopfield
network with the energy
\(
  E(\xi) = -\beta^{-1}\log\sum_a e^{\beta\,\xi^\top x_a} + \tfrac{1}{2}\xi^\top\xi
\)
has retrieval rule
$\xi^{\mathrm{new}} = X\,\mathrm{softmax}(\beta\,X^\top\xi)$, which
coincides with the transformer attention update when $X$ is the matrix
of key/value vectors. Proposition~\ref{prop:lca-hopfield} establishes
that LCA realizes this update with stored patterns $\Xi_h$ on the unit
torus and the Fourier-HRR inner product as similarity. The
unit-modulus restriction is the same one used throughout
\cite{phasor-ssm-derivation}, \S4, and preserves Plate's binding/bundling
algebra~\cite{hrr}.

\subsection{HD computing attention precedents}

VSA-based attention has been explored
in~\cite{vsa-attention,hd-binding-attention} as a binding-based
key--value associative memory. Those constructions use the
Fourier-HRR inner product for the score and complex-domain superposition
for the mix; LSA and LCA inherit those choices verbatim. The novel
elements here are (i) the head-axis restriction of the score, which
makes the operation pointwise in time, and (ii) the explicit
embedding within the three-view phasor SSM derivation.

%======================================================================
\section{Formal system summary}\label{sec:summary}
%======================================================================

\begin{definition}[Phasor LSA layer]
A Phasor Local Self-Attention layer is parameterized by three
PhasorDense projections $\Pi_Q, \Pi_K, \Pi_V$ of width $D = \nh\Dh$,
a scale $\beta\in\R_+$, and the shared carrier $\omega = 2\pi/T$. It
acts on $X\in[-1,1]^{D\times L\times B}$ via
\[
  Y_{:,l,b}
  = \tfrac{1}{\pi}\arg\!\Bigl(
      \sum_{h'}\tfrac{1}{\nh}\exp\!\bigl(\beta\,\sm(\Pi_Q(X)_{:,h,l,b},
                                                  \Pi_K(X)_{:,h',l,b})\bigr)\,
        e^{\im\pi \Pi_V(X)_{:,h',l,b}}\Bigr),
\]
broadcast over $(l, b)$. The layer admits the discrete, continuous,
and parallel forms of \eqref{eq:lsa-discrete}--\eqref{eq:lsa-parallel}
and they coincide at every lattice point $l T$
(Proposition~\ref{prop:lsa-equiv}).
\end{definition}

\begin{definition}[Phasor LCA layer]
A Phasor Local Cross-Attention layer is parameterized by an anchor
bank $Q^{*}\in[-1,1]^{D\times A}$, two PhasorDense projections
$\Pi_K, \Pi_V$ of width $D$, a head-mix tensor $M\in\R^{A\times \nh\times \nh}$
(default $\delta_{h h'}/A$), a scale $\beta\in\R_+$, and the carrier
$\omega = 2\pi/T$. It acts on $X\in[-1,1]^{D\times L\times B}$ via
\eqref{eq:lca-score}--\eqref{eq:lca-reduce}, broadcast over $(l, b)$.
It admits the three views of
\eqref{eq:lca-discrete}--\eqref{eq:lca-parallel}, which coincide at
every lattice point $l T$ (Proposition~\ref{prop:lca-equiv}). When
$V = K$ and $M$ is the default, it reduces to a modern-Hopfield
pattern retrieval on the unit torus with stored patterns $\Xi_h$
(Proposition~\ref{prop:lca-hopfield}).
\end{definition}

\medskip
\noindent\textbf{Properties.} Both LSA and LCA:
\begin{enumerate}[nosep,leftmargin=*]
  \item are pointwise in time, so they compose with the surrounding
        Phasor SSM in all three views;
  \item respect the per-channel-$\omega$ rule of
        \cite{phasor-ssm-derivation}, \S3.5---all heads share $\omega$,
        only $\lambda$ varies per channel;
  \item realize the Fourier-HRR algebra on a per-slice basis: the
        score is HRR similarity, the mix is HRR bundling, the output
        $\arg$ is the HRR phase-extraction nonlinearity.
\end{enumerate}

The companion document
\texttt{docs/local\_attention\_companion.md} traces these definitions
to the existing PhasorNetworks.jl source tree and lays out an
implementation and experimental test plan.

%======================================================================
\begin{thebibliography}{99}

\bibitem{phasor-ssm-derivation}
W.\ Olin-Ammentorp.
\newblock The Phasor State-Space Model: A Derivation from
  Resonate-and-Fire Dynamics to Holographic Distributed Representations.
\newblock Companion document, PhasorNetworks.jl
  (\texttt{docs/phasor\_ssm\_derivation.tex}).

\bibitem{hopfield-is-all-you-need}
H.\ Ramsauer, B.\ Sch{\"a}fl, J.\ Lehner, P.\ Seidl, M.\ Widrich,
T.\ Adler, L.\ Gruber, M.\ Holzleitner, M.\ Pavlovi{\'c},
G.\ Kraml Sandner, et al.
\newblock Hopfield Networks Is All You Need.
\newblock In \emph{International Conference on Learning Representations
  (ICLR)}, 2021.

\bibitem{longformer}
I.\ Beltagy, M.\ E.\ Peters, and A.\ Cohan.
\newblock Longformer: The Long-Document Transformer.
\newblock \emph{arXiv preprint} arXiv:2004.05150, 2020.

\bibitem{longformer-followups}
J.\ Ainslie, S.\ Onta{\~n}{\'o}n, C.\ Alberti, V.\ Cvicek, Z.\ Fisher,
P.\ Pham, A.\ Ravula, S.\ Sanghai, Q.\ Wang, and L.\ Yang.
\newblock ETC: Encoding Long and Structured Inputs in Transformers.
\newblock In \emph{Empirical Methods in Natural Language Processing
  (EMNLP)}, 2020.

\bibitem{efficient-transformers}
Y.\ Tay, M.\ Dehghani, D.\ Bahri, and D.\ Metzler.
\newblock Efficient Transformers: A Survey.
\newblock \emph{ACM Computing Surveys}, 55(6), 2022.

\bibitem{mamba}
A.\ Gu and T.\ Dao.
\newblock Mamba: Linear-Time Sequence Modeling with Selective State
  Spaces.
\newblock In \emph{Conference on Language Modeling (COLM)}, 2024.

\bibitem{retnet}
Y.\ Sun, L.\ Dong, S.\ Huang, S.\ Ma, Y.\ Xia, J.\ Xue, J.\ Wang, and
F.\ Wei.
\newblock Retentive Network: A Successor to Transformer for Large
  Language Models.
\newblock \emph{arXiv preprint} arXiv:2307.08621, 2023.

\bibitem{s5}
J.\ T.\ H.\ Smith, A.\ Warrington, and S.\ W.\ Linderman.
\newblock Simplified State Space Layers for Sequence Modeling.
\newblock In \emph{International Conference on Learning Representations
  (ICLR)}, 2023.

\bibitem{hippo}
A.\ Gu, T.\ Dao, S.\ Ermon, A.\ Rudra, and C.\ R{\'e}.
\newblock HiPPO: Recurrent Memory with Optimal Polynomial Projections.
\newblock In \emph{Neural Information Processing Systems (NeurIPS)},
  2020.

\bibitem{hrr}
T.\ A.\ Plate.
\newblock Holographic Reduced Representations.
\newblock \emph{IEEE Transactions on Neural Networks}, 6(3):623--641,
  1995.

\bibitem{vsa-attention}
E.\ P.\ Frady, D.\ Kleyko, and F.\ T.\ Sommer.
\newblock A Theory of Sequence Indexing and Working Memory in
  Recurrent Neural Networks.
\newblock \emph{Neural Computation}, 30(6):1449--1513, 2018.

\bibitem{hd-binding-attention}
P.\ Neubert, S.\ Schubert, K.\ Schlegel, and P.\ Protzel.
\newblock Vector Semantic Representations as Descriptors for Visual
  Place Recognition.
\newblock In \emph{Robotics: Science and Systems}, 2021.

\end{thebibliography}

\end{document}
